% ============================================================
%  PRÉSENTATION BEAMER — YouTube Quality Analyzer
%  Système Multi-Agents LLM pour l'Analyse de la Qualité
%  des Commentaires YouTube à Vocation Pédagogique
% ============================================================
% Compilation : pdflatex presentation.tex (2 passes)
% ============================================================

\documentclass[aspectratio=169, 12pt]{beamer}

% ── Thème ────────────────────────────────────────────────────
\usetheme{Madrid}
\usecolortheme{whale}
\usefonttheme{structurebold}

% ── Encodage ─────────────────────────────────────────────────
\usepackage[utf8]{inputenc}
\usepackage[T1]{fontenc}
\usepackage[french]{babel}
\usepackage{pifont}
\newcommand{\cmark}{\textcolor{accentgreen}{\ding{51}}}
\newcommand{\xmark}{\textcolor{red}{\ding{55}}}

% ── Packages ─────────────────────────────────────────────────
\usepackage{graphicx}
\usepackage{booktabs}
\usepackage{array}
\usepackage{xcolor}
\usepackage{amsmath}
\usepackage{tikz}
\usetikzlibrary{positioning, arrows.meta, shapes.geometric, fit, backgrounds}
\usepackage{pgfplots}
\pgfplotsset{compat=1.18}
\usepackage{listings}
\usepackage{tcolorbox}
\tcbuselibrary{skins}

% ── Couleurs personnalisées ───────────────────────────────────
\definecolor{ytred}{RGB}{255, 0, 0}
\definecolor{mainblue}{RGB}{0, 60, 120}
\definecolor{accentgreen}{RGB}{40, 160, 40}
\definecolor{accentorange}{RGB}{220, 120, 0}
\definecolor{lightgray}{RGB}{245, 245, 245}

\setbeamercolor{palette primary}{bg=mainblue, fg=white}
\setbeamercolor{palette secondary}{bg=mainblue!80, fg=white}
\setbeamercolor{palette tertiary}{bg=mainblue!60, fg=white}
\setbeamercolor{structure}{fg=mainblue}
\setbeamercolor{block title}{bg=mainblue, fg=white}
\setbeamercolor{block body}{bg=lightgray}
\setbeamercolor{alertblock title}{bg=ytred!80!black, fg=white}

% ── Listings (code) ──────────────────────────────────────────
\lstset{
  basicstyle=\ttfamily\tiny,
  breaklines=true,
  backgroundcolor=\color{lightgray},
  frame=single,
  showstringspaces=false,
}

% ── Navigation symbols ───────────────────────────────────────
\setbeamertemplate{navigation symbols}{}
\setbeamertemplate{footline}{%
  \begin{beamercolorbox}[wd=\paperwidth, ht=2.5ex, dp=1ex,
                         leftskip=0.4cm, rightskip=0.4cm]{palette primary}%
    \usebeamerfont{author in head/foot}%
    \insertshortauthor\hfill
    \insertshorttitle\hfill
    \insertframenumber\,/\,\inserttotalframenumber%
  \end{beamercolorbox}%
}

% ── Informations ─────────────────────────────────────────────
\title[YouTube Quality Analyzer]{YouTube Quality Analyzer}
\subtitle{Système Multi-Agents LLM pour l'Analyse\\
          de la Qualité des Commentaires Pédagogiques}
\author[MASSO YETNA · MOUKODI · WATAT]{
  \textbf{MASSO YETNA} Ferdinand Hervé \\[0.2em]
  \textbf{MOUKODI} Claudia Lauren Julienne \\[0.2em]
  \textbf{WATAT YONDEP} Stive Kevin
}
\institute{}
\date{Avril 2026}

% ═══════════════════════════════════════════════════════════════
\begin{document}
% ═══════════════════════════════════════════════════════════════

% ── SLIDE 1 — Page de titre ──────────────────────────────────
\begin{frame}[plain]

  %% ── Bandeau titre ────────────────────────────────────────
  \begin{beamercolorbox}[wd=\paperwidth, ht=3.0cm, dp=0.4cm, center]{palette primary}
    \vspace{0.45cm}
    {\LARGE\textbf{YouTube Quality Analyzer}}\\[0.3em]
    {\small\itshape Système Multi-Agents LLM pour l'Analyse\\
    de la Qualité des Commentaires Pédagogiques}
  \end{beamercolorbox}

  \vspace{0.35cm}

  %% ── Auteurs + framework ──────────────────────────────────
  \begin{columns}[T]
    \begin{column}{0.54\textwidth}
      \begin{center}
        \small
        \textbf{MASSO YETNA} Ferdinand Hervé\\[0.18em]
        \textbf{MOUKODI} Claudia Lauren Julienne\\[0.18em]
        \textbf{WATAT YONDEP} Stive Kevin\\[0.35em]
        \textcolor{gray}{\textit{Avril 2026}}
      \end{center}
    \end{column}
    \begin{column}{0.43\textwidth}
      \begin{center}
        \scriptsize\textcolor{gray}{\textit{Projet de recherche appliquée\\
        en intelligence artificielle pédagogique}}\\[0.35em]
        \scriptsize\textbf{Framework :}\\[0.1em]
        \scriptsize LangGraph · FastAPI\\
        GPT-4o-mini · YouTube Data API v3
      \end{center}
    \end{column}
  \end{columns}

  \vspace{0.25cm}
  {\color{mainblue!30}\rule{\linewidth}{0.4pt}}
  \vspace{0.15cm}

  %% ── Logos ────────────────────────────────────────────────
  \begin{center}
    \begin{tabular}{@{\hspace{0.3cm}}c@{\hspace{0.5cm}}c@{\hspace{0.5cm}}c@{\hspace{0.5cm}}c@{\hspace{0.5cm}}c@{\hspace{0.5cm}}c@{\hspace{0.3cm}}}
      \IfFileExists{logo_cnrs.png}{\includegraphics[height=0.85cm]{logo_cnrs.png}}{} &
      \IfFileExists{logo_uca.png}{\includegraphics[height=0.85cm]{logo_uca.png}}{} &
      \IfFileExists{logo_ummisco.png}{\includegraphics[height=0.85cm]{logo_ummisco.png}}{} &
      \IfFileExists{logo_inp.png}{\includegraphics[height=0.85cm]{logo_inp.png}}{} &
      \IfFileExists{logo_enspy.png}{\includegraphics[height=0.85cm]{logo_enspy.png}}{} &
      \IfFileExists{logo_uy1.png}{\includegraphics[height=0.85cm]{logo_uy1.png}}{}
    \end{tabular}
  \end{center}

\end{frame}

% ── SLIDE 2 — Plan ───────────────────────────────────────────
\begin{frame}{Plan de la présentation}
  \tableofcontents[hideallsubsections]
\end{frame}

\section{Le Problème}

% ── SLIDE 3 — Problématique ──────────────────────────────────
\begin{frame}{Problématique}

  \begin{block}{Contexte}
    \begin{itemize}
      \item Des millions de vidéos pédagogiques sur YouTube
      \item La qualité perçue est \textbf{difficile à mesurer automatiquement}
      \item Les commentaires = signal riche, massif, non structuré
    \end{itemize}
  \end{block}

  \vspace{0.4cm}
  \begin{alertblock}{Question de recherche}
    \centering
    \textit{Un système d'intelligence artificielle peut-il lire\\
    les commentaires YouTube et en déduire si une vidéo\\
    vaut la peine d'être regardée pour apprendre ?}
  \end{alertblock}

\end{frame}

% ── SLIDE — Persona 1 ────────────────────────────────────────
\begin{frame}{Persona 1 — L'Étudiante}
  \begin{columns}[T]
    \begin{column}{0.40\textwidth}
      \begin{center}
        \vspace{0.5cm}
        {\LARGE\textcolor{mainblue}{\textbf{Aminata}}}\\[0.3em]
        {\large 20 ans · Étudiante}\\[0.2em]
        \textit{Licence Sciences Économiques}
      \end{center}
    \end{column}
    \begin{column}{0.57\textwidth}
      \begin{block}{Contexte}
        \small Utilise YouTube pour compléter ses cours universitaires.
        Face à des dizaines de vidéos sur le même sujet, elle ne sait
        pas laquelle choisir.
      \end{block}
      \vspace{0.2cm}
      \begin{alertblock}{Frustration}
        \small Les likes et les vues ne garantissent pas la qualité
        pédagogique réelle d'une vidéo.
      \end{alertblock}
      \vspace{0.2cm}
      \begin{exampleblock}{Ce que la solution lui apporte}
        \small Un score de qualité perçue agrégé lui permet de choisir
        la meilleure vidéo \textbf{en quelques secondes}, sans lire
        des dizaines de commentaires.
      \end{exampleblock}
    \end{column}
  \end{columns}
\end{frame}

% ── SLIDE — Persona 2 ────────────────────────────────────────
\begin{frame}{Persona 2 — L'Apprenant Sourd}
  \begin{columns}[T]
    \begin{column}{0.40\textwidth}
      \begin{center}
        \vspace{0.5cm}
        {\LARGE\textcolor{mainblue}{\textbf{Théo}}}\\[0.3em]
        {\large 22 ans · Étudiant}\\[0.2em]
        \textit{Licence Informatique\\sourd de naissance}
      \end{center}
    \end{column}
    \begin{column}{0.57\textwidth}
      \begin{block}{Contexte}
        \small S'appuie exclusivement sur les sous-titres et les visuels
        pour apprendre via YouTube. Il ne peut pas évaluer la clarté
        audio d'une vidéo.
      \end{block}
      \vspace{0.2cm}
      \begin{alertblock}{Frustration}
        \small Aucune métrique YouTube ne reflète la qualité des
        sous-titres ni l'accessibilité réelle du contenu.
      \end{alertblock}
      \vspace{0.2cm}
      \begin{exampleblock}{Ce que la solution lui apporte}
        \small Détection des signaux de clarté et de compréhensibilité
        dans les commentaires — indicateur indirect d'accessibilité.
      \end{exampleblock}
    \end{column}
  \end{columns}
\end{frame}

\section{Notre Solution}

% ── SLIDE — Extension Chrome (1/4) ───────────────────────────
\begin{frame}{Extension Chrome — Présentation (1/4)}

  \begin{block}{Qu'est-ce que c'est ?}
    \normalsize
    Une extension navigateur \textbf{Chrome MV3} qui s'intègre directement
    dans l'interface YouTube. L'utilisateur n'a pas besoin de quitter la
    page vidéo pour obtenir une analyse de qualité des commentaires.
  \end{block}

  \vspace{0.4cm}

  \begin{block}{Ce qu'elle fait}
    \begin{itemize}
      \item Détecte automatiquement l'URL de la vidéo YouTube en cours de visionnage
      \item Envoie la requête d'analyse vers l'API backend (FastAPI)
      \item Affiche le \textbf{score de qualité perçue} directement en overlay sur la page
      \item Présente un résumé sentimental et les thèmes dominants en temps réel
    \end{itemize}
  \end{block}

\end{frame}

% ── SLIDE — Extension Chrome (2/4) ───────────────────────────
\begin{frame}{Extension Chrome — Ce qui est fonctionnel (2/4)}

  \begin{columns}[T]
    \begin{column}{0.48\textwidth}
      \begin{block}{Fonctionnalités de base}
        \begin{itemize}
          \item Injection de contenu dans la page YouTube
          \item Communication avec l'API via requêtes HTTP
          \item Interface d'affichage du score de qualité intégrée
        \end{itemize}
      \end{block}
    \end{column}
    \begin{column}{0.48\textwidth}
      \begin{exampleblock}{Module Q\&A intégré}
        \small
        \begin{itemize}
          \item Pose une question sur la vidéo en cours
          \item Réponse générée à partir des commentaires et de la transcription
          \item Résultat affiché directement dans l'overlay
        \end{itemize}
      \end{exampleblock}
    \end{column}
  \end{columns}

\end{frame}

% ── SLIDE — Extension Chrome (3/4) ───────────────────────────
\begin{frame}{Extension Chrome — Système de quiz (3/4)}

  \begin{block}{Qu'est-ce que le système de quiz ?}
    \normalsize
    Un module intégré à l'extension qui génère automatiquement des
    \textbf{QCM} à partir du contenu de la vidéo YouTube en cours de
    visionnage, sans quitter la page.
  \end{block}

  \vspace{0.4cm}

  \begin{exampleblock}{Fonctionnement}
    \begin{itemize}
      \item Génération de questions à choix multiples à partir de la transcription et des commentaires
      \item Chaque question est \textbf{horodatée} avec une citation source obligatoire
      \item Correction immédiate — l'apprenant voit sa réponse validée ou corrigée en temps réel
      \item Permet de tester la compréhension \textbf{directement dans l'interface YouTube}
    \end{itemize}
  \end{exampleblock}

\end{frame}

% ── SLIDE — Extension Chrome (4/4) ───────────────────────────
\begin{frame}{Extension Chrome — Garantie qualité (4/4)}

  \begin{alertblock}{Garantie anti-hallucination}
    \normalsize
    Chaque question est ancrée sur une \textbf{citation horodatée} extraite
    du contenu réel de la vidéo — le LLM ne peut pas inventer de réponse
    hors source.
  \end{alertblock}

  \vspace{0.5cm}

  \begin{block}{Pourquoi c'est important ?}
    \begin{itemize}
      \item Les réponses générées restent \textbf{vérifiables} par l'apprenant
      \item La citation horodatée permet de \textbf{retrouver le passage} dans la vidéo
      \item Aucune information inventée ne peut passer dans le quiz
      \item Renforce la \textbf{confiance} de l'utilisateur dans le système
    \end{itemize}
  \end{block}

\end{frame}

% ── SLIDE — API REST (1/2) ────────────────────────────────────
\begin{frame}{API REST — Architecture \& Utilité (1/2)}

  \begin{block}{Architecture FastAPI}
    \begin{itemize}
      \item Endpoint principal \texttt{POST /analyze} — déclenche le pipeline complet
      \item \textbf{Streaming SSE (Server-Sent Events)} — les résultats arrivent progressivement
      \item Réponse initiale en \textbf{$\approx$8 secondes} sur 300 commentaires
      \item Enrichissement du corpus complet en arrière-plan (thread daemon)
    \end{itemize}
  \end{block}

  \vspace{0.4cm}

  \begin{block}{Utilité}
    \normalsize
    Tout client — extension Chrome, application web, script Python — peut
    appeler l'API pour obtenir un \textbf{score de qualité perçue} sur
    n'importe quelle vidéo YouTube pédagogique, sans connaître les détails
    du pipeline interne.
  \end{block}

\end{frame}

% ── SLIDE — API REST (2/2) ────────────────────────────────────
\begin{frame}[fragile]{API REST — Flux temps réel (2/2)}

  \begin{block}{Comment fonctionne le Streaming SSE ?}
    \normalsize
    Le client reçoit les événements au fur et à mesure que chaque agent
    termine son traitement — sans attendre la fin du pipeline complet.
  \end{block}

  \vspace{0.3cm}

  \begin{exampleblock}{Exemple de flux SSE reçu par le client}
    \lstset{basicstyle=\ttfamily\small}
    \begin{lstlisting}
t+0s   event: started       -> pipeline lancé
t+2s   event: collected     -> commentaires récupérés
t+3s   event: preprocessed  -> nettoyage terminé
t+6s   event: scores        -> scores agents disponibles
t+7s   event: global        -> score global calculé
t+8s   event: final         -> score final avec pertinence
t+8s   event: done          -> pipeline terminé
    \end{lstlisting}
  \end{exampleblock}

\end{frame}

% ── SLIDE — Collecte de données (1/2) ────────────────────────
\begin{frame}{Projet de collecte de données — Présentation (1/2)}

  \begin{block}{Un projet open source dédié}
    \normalsize
    Avant de construire le système d'analyse, nous avons conçu un
    \textbf{pipeline de collecte automatique} de commentaires YouTube
    via la \textbf{YouTube Data API v3}. Ce projet est indépendant
    et réutilisable par la communauté.\\[0.5em]
    \textbf{Code source disponible :}\\[0.2em]
    \href{https://github.com/FerdinandMY/data_collection_youtube}
      {\texttt{github.com/FerdinandMY/data\_collection\_youtube}}
  \end{block}

  \vspace{0.4cm}

  \begin{exampleblock}{Corpus produit}
    \centering
    \begin{tabular}{lc}
      \toprule
      \textbf{Métrique} & \textbf{Valeur} \\
      \midrule
      Commentaires collectés & \textbf{42\,860} \\
      Vidéos pédagogiques & \textbf{591} \\
      Langues & FR 65\% · EN 35\% \\
      Licence & MIT (open source) \\
      \bottomrule
    \end{tabular}
  \end{exampleblock}

\end{frame}

% ── SLIDE — Collecte de données (2/2) ────────────────────────
\begin{frame}{Projet de collecte de données — Fonctionnement (2/2)}

  \begin{block}{Comment le pipeline collecte les données ?}
    \begin{itemize}
      \item Recherche automatique de vidéos par \textbf{mots-clés pédagogiques}
      \item Filtrage par langue, durée minimale et niveau d'engagement
      \item Récupération des commentaires via pagination de l'API YouTube
      \item Export structuré en \textbf{CSV} (métadonnées vidéo + commentaires)
      \item Gestion automatique des \textbf{quotas API} (10\,000 unités/jour)
    \end{itemize}
  \end{block}

  \vspace{0.4cm}

  \begin{block}{Pourquoi ce projet est important ?}
    \normalsize
    Il constitue la \textbf{première base de données publique} de commentaires
    YouTube pédagogiques en français, directement utilisable pour entraîner
    ou évaluer d'autres systèmes d'analyse de qualité.
  \end{block}

\end{frame}

\section{Architecture du Système}

% ── SLIDE 4a — Architecture Pipeline ─────────────────────────
\begin{frame}{Architecture du Système — Pipeline d'Analyse}
\begin{center}
  \IfFileExists{architecture_pipeline.png}{%
    \includegraphics[width=0.92\textwidth, height=0.82\textheight,
                     keepaspectratio]{architecture_pipeline.png}%
  }{%
    \textcolor{gray}{\textit{[Image non trouvée : placer \texttt{architecture\_pipeline.png}
    dans le dossier du projet]}}%
  }
\end{center}
\end{frame}

% ── SLIDE 4b — Outils LangChain ───────────────────────────────
\begin{frame}{Architecture du Système — Outils \& Couches Transverses}
\begin{center}
  \IfFileExists{architecture_tools.png}{%
    \includegraphics[width=0.92\textwidth, height=0.82\textheight,
                     keepaspectratio]{architecture_tools.png}%
  }{%
    \textcolor{gray}{\textit{[Image non trouvée : placer \texttt{architecture\_tools.png}
    dans le dossier du projet]}}%
  }
\end{center}
\end{frame}

% ── SLIDE 6 — IA Générative : Prompts ────────────────────────
\begin{frame}{IA Générative — Ingénierie des prompts}

  \begin{columns}[T]
    \begin{column}{0.48\textwidth}
      \begin{block}{Structure de chaque prompt}
        \begin{enumerate}
          \item \textbf{Rôle} — identité de l'agent
          \item \textbf{Contexte} — données à analyser
          \item \textbf{Raisonnement} — CoT / ToT / SC
          \item \textbf{Format} — JSON strict
        \end{enumerate}
      \end{block}

      \vspace{0.1cm}
      \begin{alertblock}{Anti-hallucination}
        \small
        \begin{itemize}\setlength{\itemsep}{1pt}
          \item Q\&A : répondre \textbf{UNIQUEMENT} sur les sources
          \item QCM : citation horodatée \textbf{obligatoire}
          \item A7 : vote majoritaire (3 passes)
        \end{itemize}
      \end{alertblock}
    \end{column}

    \begin{column}{0.48\textwidth}
      \begin{block}{Techniques de raisonnement}
        \small
        \textbf{Chain-of-Thought (A3)} \\
        \textit{Thought 1 : VADER → Thought 2 : marqueurs → Conclusion}
        \vspace{0.15cm}

        \textbf{Tree of Thoughts (A4)} \\
        \textit{Br. 1 : profondeur · Br. 2 : lexique · Br. 3 : épistémique}
        \vspace{0.15cm}

        \textbf{Self-Consistency (A7)} \\
        \textit{3 évaluations parallèles $\rightarrow$ vote majoritaire}
      \end{block}
    \end{column}
  \end{columns}
\end{frame}

% ── SLIDE — Techniques de raisonnement ──────────────────────
\begin{frame}{Techniques de raisonnement — ReAct, CoT, ToT}

  \begin{columns}[T]

    \begin{column}{0.32\textwidth}
      \begin{block}{ReAct — Agent A3}
        \small\textbf{Sentiment}\\[0.3em]
        \scriptsize
        Cycle itératif :\\
        \textbf{Thought} → \textbf{Action} → \textbf{Obs.} → \textbf{Result}\\[0.4em]
        Étapes réelles :
        \begin{enumerate}\setlength{\itemsep}{1pt}
          \item Thought 1 — score VADER
          \item Thought 2 — correction sarcastique
          \item Thought 3 — tonalité finale
          \item Result — JSON
        \end{enumerate}
        \vspace{0.2cm}
        \textcolor{gray}{\tiny Tools appelés \textbf{avant} le LLM :\\
        \texttt{vader\_sentiment} · \texttt{detect\_sarcasm\_markers}}
      \end{block}
    \end{column}

    \begin{column}{0.32\textwidth}
      \begin{block}{CoT + ToT conditionnel — A4}
        \small\textbf{Discours}\\[0.3em]
        \scriptsize
        \textbf{CoT} (toujours) :\\
        Thought 1 info → Thought 2 arg.\\
        → Thought 3 construct. → Result\\[0.4em]
        \textbf{ToT} si score ambigu $[0.35, 0.65]$ :\\
        \begin{itemize}\setlength{\itemsep}{1pt}
          \item Br. 1 — profondeur du contenu
          \item Br. 2 — qualité de la discussion
          \item Br. 3 — valeur épistémique
          \item Synthèse des 3 branches
        \end{itemize}
        \textcolor{gray}{\tiny ToT remplace CoT si non-fallback}
      \end{block}
    \end{column}

    \begin{column}{0.32\textwidth}
      \begin{block}{ToT + Self-Consistency — A7}
        \small\textbf{Topic Matcher}\\[0.3em]
        \scriptsize
        \textbf{ToT systématique} (3 branches) :\\
        \begin{itemize}\setlength{\itemsep}{1pt}
          \item Br. 1 — alignement contenu
          \item Br. 2 — niveau audience
          \item Br. 3 — valeur pédagogique
        \end{itemize}
        \vspace{0.2cm}
        \textbf{Self-Consistency} (3 runs $\parallel$) :\\
        \texttt{ThreadPoolExecutor} · vote\\
        majoritaire label + moyenne score\\
        Si consensus $< 2/3$ → flag\\
        \texttt{low\_consensus} dans rapport
      \end{block}
    \end{column}

  \end{columns}

  \vspace{0.2cm}
  \begin{center}
    \scriptsize
    \textcolor{gray}{A5 utilise un \textbf{CoT léger} uniquement si confiance heuristique $< 0.75$ —
    sinon résultat heuristique direct sans LLM.}
  \end{center}

\end{frame}

% ── SLIDE 7 — Tools & Function Calling ───────────────────────
\begin{frame}{Tools \& Function Calling}

  \begin{columns}[T]
    \begin{column}{0.44\textwidth}
      \begin{center}
      \begin{tikzpicture}[
        box/.style={rectangle, rounded corners=3pt, draw=mainblue,
                    fill=mainblue!10, minimum width=2.4cm, minimum height=0.55cm,
                    font=\scriptsize, text centered},
        tool/.style={rectangle, rounded corners=3pt, draw=accentorange,
                     fill=accentorange!10, minimum width=2.4cm, minimum height=0.55cm,
                     font=\scriptsize, text centered},
        arr/.style={-{Stealth[length=3pt]}, thick},
        node distance=0.35cm and 0.8cm,
      ]
        \node[box] (p) {Prompt + données};
        \node[box, below=0.5cm of p] (llm) {LLM GPT-4o-mini};
        \node[tool, right=0.9cm of llm] (t) {Tool Python};
        \node[box, below=0.5cm of llm] (r) {Résultat JSON};
        \draw[arr] (p) -- (llm) node[midway,right,font=\tiny]{invoke};
        \draw[arr] (llm) -- (t) node[midway,above,font=\tiny]{tool\_call};
        \draw[arr] (t) -- (r) node[midway,right,font=\tiny]{résultat};
        \draw[arr] (llm) -- (r);
      \end{tikzpicture}
      \end{center}
    \end{column}

    \begin{column}{0.54\textwidth}
      \scriptsize
      \begin{tabular}{@{}lp{2.9cm}@{}}
        \toprule
        \textbf{Tool} & \textbf{Rôle} \\
        \midrule
        \texttt{compute\_vader} & Scores VADER — ancrage A3 \\
        \texttt{compute\_text\_stats} & Diversité lexicale (A4) \\
        \texttt{get\_top\_liked} & Top commentaires par likes (A4) \\
        \texttt{svm\_spam\_detector} & Heuristique spam (A5) \\
        \texttt{count\_repeated\_chars} & Détection bots (A5) \\
        \texttt{retrieve\_context} & Transcription + top comments (Q\&A) \\
        \bottomrule
      \end{tabular}

      \vspace{0.15cm}
      \begin{block}{\small Avantage}
        \scriptsize LLM = raisonnement · Tool = calcul déterministe\\
        $\rightarrow$ fallback si LLM indisponible
      \end{block}
    \end{column}
  \end{columns}
\end{frame}

% (suite Notre Solution)

\section{Réalisations}

% ── SLIDE — Formules de scoring ──────────────────────────────
\begin{frame}{Formules de scoring}

  \begin{block}{Score Global — 3 dimensions pondérées}
    \begin{equation*}
      \text{Score\_Global} = \underbrace{0{,}35}_{\text{Sentiment}} \times S
                           + \underbrace{0{,}40}_{\text{Discours}} \times D
                           + \underbrace{0{,}25}_{\text{Bruit}} \times N
    \end{equation*}
  \end{block}

  \vspace{0.3cm}

  \begin{block}{Score Final — pertinence thématique}
    \begin{equation*}
      \text{Score\_Final} = 0{,}60 \times \text{Score\_Global}
                          + 0{,}40 \times \text{Score\_Pertinence}
    \end{equation*}
  \end{block}

  \vspace{0.3cm}
  \begin{alertblock}{Interprétation}
    \small
    \begin{itemize}
      \item $S$ = score sentiment · $D$ = score discours · $N$ = score bruit inversé
      \item \texttt{Score\_Pertinence} calculé par A7 (Topic Matcher) via Self-Consistency
      \item Score Final $\in [0, 100]$ — seuil de qualité pédagogique : $\geq 60$
    \end{itemize}
  \end{alertblock}

\end{frame}

% ── SLIDE — Justification des poids ──────────────────────────
\begin{frame}{Justification des poids — Micro-ablation}

  \begin{columns}[T]
    \begin{column}{0.54\textwidth}
      \begin{block}{Étude d'ablation sur les pondérations}
        \centering\small
        \begin{tabular}{lc}
          \toprule
          Configuration & Pearson $r$ \\
          \midrule
          Uniforme (0,33 / 0,33 / 0,33) & 0,8421 \\
          Sentiment dominant & 0,8104 \\
          \textbf{Discours dominant (retenu)} & \textbf{0,8867} \\
          Bruit dominant & 0,7932 \\
          \bottomrule
        \end{tabular}
      \end{block}
    \end{column}
    \begin{column}{0.43\textwidth}
      \begin{alertblock}{Pourquoi $w_D = 0{,}40$ ?}
        \small La qualité du discours — connecteurs logiques,
        diversité lexicale, marqueurs argumentatifs — est le
        \textbf{meilleur prédicteur empirique} de la valeur
        pédagogique perçue.
      \end{alertblock}
    \end{column}
  \end{columns}

\end{frame}

\section{Éthique}

% ── SLIDE — Éthique (1/2) ────────────────────────────────────
\begin{frame}{Considérations éthiques (1/2) — Données \& Vie privée}

  \begin{block}{Données collectées}
    \begin{itemize}
      \item Commentaires \textbf{publics} uniquement — aucune donnée privée
      \item Collecte via YouTube Data API v3 dans le respect des
            \textbf{conditions d'utilisation} Google
      \item Aucun identifiant personnel stocké — les auteurs de
            commentaires ne sont pas tracés
      \item Dataset publié sous licence \textbf{MIT} (open source)
    \end{itemize}
  \end{block}

  \vspace{0.4cm}
  \begin{exampleblock}{Mesures prises pour limiter les biais}
    \small
    Annotation croisée inter-modèles (GPT-4o-mini $\times$ Llama-3.1-70B,
    $\kappa = 0{,}779$) pour limiter le biais auto-référentiel — niveau 3
    (annotation humaine) identifié comme perspective prioritaire.\\[0.3em]
    Corpus majoritairement FR/EN — résultats non généralisables à d'autres
    langues sans adaptation spécifique.
  \end{exampleblock}

\end{frame}

% ── SLIDE — Éthique (2/2) ────────────────────────────────────
\begin{frame}{Considérations éthiques (2/2) — IA \& Usage pédagogique}

  \begin{columns}[T]
    \begin{column}{0.48\textwidth}
      \begin{block}{Transparence du système}
        \begin{itemize}
          \item Le score est \textbf{explicable} : chaque dimension
                (sentiment, discours, bruit) est visible dans le rapport
          \item Les citations horodatées du quiz permettent à l'apprenant
                de \textbf{vérifier} chaque réponse générée
          \item L'utilisateur est informé que le score est basé sur
                les \textbf{commentaires}, pas sur le contenu vidéo direct
        \end{itemize}
      \end{block}
    \end{column}
    \begin{column}{0.48\textwidth}
      \begin{alertblock}{Risques identifiés}
        \begin{itemize}
          \item Un score bas peut \textbf{décourager} l'accès à une vidéo
                de qualité sous-commentée
          \item Le système ne remplace pas le \textbf{jugement humain}
                d'un enseignant
          \item Hallucinations résiduelles possibles malgré les
                5 niveaux de garde
        \end{itemize}
      \end{alertblock}
    \end{column}
  \end{columns}

  \vspace{0.3cm}
  \begin{block}{Engagement}
    \small
    Le système est conçu comme un \textbf{outil d'aide à la décision},
    non comme un arbitre unique. L'apprenant reste acteur de son parcours —
    le score est un signal parmi d'autres, pas une vérité absolue.
  \end{block}

\end{frame}

% ── SLIDE — Conclusion ───────────────────────────────────────
\begin{frame}{Conclusion}

  \begin{block}{Bilan}
    \begin{itemize}
      \item Pipeline multi-agents \textbf{opérationnel} et évalué sur corpus réel
      \item 4 hypothèses validées avec \textbf{marges significatives}
            (Pearson $r = 0{,}8867$, F1-macro $= 0{,}9652$, fallback $= 0\,\%$)
      \item Modules pédagogiques Q\&A et QCM \textbf{intégrés} à l'extension Chrome
      \item Architecture \textbf{extensible} et open source (MIT)
    \end{itemize}
  \end{block}

  \vspace{0.4cm}

  \begin{center}
    {\Large\textbf{\textcolor{mainblue}{Merci !}}}\\[0.4em]
    \rule{0.5\linewidth}{0.4pt}\\[0.3em]
    \small
    \textbf{Contact :}
    MASSO YETNA Ferdinand Hervé ·
    MOUKODI Claudia Lauren Julienne ·
    WATAT YONDEP Stive Kevin\\[0.2em]
    \scriptsize\textit{ENSPY / Université de Yaoundé I — Avril 2026}
  \end{center}

\end{frame}

% ── SLIDE — Perspectives ─────────────────────────────────────
\begin{frame}{Perspectives \& Pistes d'amélioration}

  \begin{columns}[T]
    \begin{column}{0.48\textwidth}
      \begin{block}{Amélioration du système}
        \begin{itemize}\setlength{\itemsep}{3pt}
          \item Annotation humaine (niveau 3) — réduire le biais
                auto-référentiel LLM
          \item Extension corpus au-delà de 100 commentaires annotés
                (contrainte GPU T4)
          \item Améliorer le rappel A5 : 12\,\% $\to$ 40\,\%+
                pour détecter le bruit subtil
          \item Déploiement cloud stable (Railway / Render)
        \end{itemize}
      \end{block}
    \end{column}
    \begin{column}{0.48\textwidth}
      \begin{exampleblock}{Généralisation \& Accessibilité}
        \begin{itemize}\setlength{\itemsep}{3pt}
          \item Généralisation au-delà du domaine STEM
                (langues, sciences humaines)
          \item Accessibilité : détection explicite des signaux
                sous-titres pour les apprenants sourds
          \item Extension Chrome MV3 complète
                + publication Chrome Web Store
          \item \textcolor{accentgreen}{\textbf{Acquis :}} annotation
                croisée niveau 2 ($\kappa=0{,}779$) \cmark
        \end{itemize}
      \end{exampleblock}
    \end{column}
  \end{columns}

\end{frame}

% ── SLIDE 16 — Références ────────────────────────────────────
\section*{Références}
\begin{frame}[allowframebreaks]{Références bibliographiques}
  \footnotesize
  \setlength{\parskip}{4pt}
  \begin{thebibliography}{99}

    \bibitem{hutto2014}
      \textbf{Hutto, C. \& Gilbert, E.} (2014).
      VADER: A Parsimonious Rule-based Model for Sentiment Analysis of
      Social Media Text.
      \textit{Proc. ICWSM}, 8, 216--225.

    \bibitem{wei2022}
      \textbf{Wei, J. et al.} (2022).
      Chain-of-Thought Prompting Elicits Reasoning in Large Language Models.
      \textit{NeurIPS 2022}.

    \bibitem{yao2023}
      \textbf{Yao, S. et al.} (2023).
      Tree of Thoughts: Deliberate Problem Solving with Large Language Models.
      \textit{NeurIPS 2023}.

    \bibitem{wang2023}
      \textbf{Wang, X. et al.} (2023).
      Self-Consistency Improves Chain of Thought Reasoning in LLMs.
      \textit{ICLR 2023}.

    \bibitem{park2016}
      \textbf{Park, S.} (2016).
      Learning Effectiveness of YouTube Videos as Online Learning Resources.
      \textit{Interdisciplinary Journal of E-Learning \& Learning Objects},
      12, 327--345.

    \bibitem{cohen1960}
      \textbf{Cohen, J.} (1960).
      A coefficient of agreement for nominal scales.
      \textit{Educational and Psychological Measurement}, 20(1), 37--46.

    \bibitem{langgraph2024}
      \textbf{LangChain AI.} (2024).
      LangGraph: Building Stateful, Multi-Actor Applications with LLMs.
      \texttt{https://github.com/langchain-ai/langgraph}.

    \bibitem{openai2024}
      \textbf{OpenAI.} (2024).
      GPT-4o Technical Report.
      \texttt{https://openai.com/research/gpt-4o}.

  \end{thebibliography}
\end{frame}

% ═══════════════════════════════════════════════════════════════
\end{document}
